\documentclass[12pt,a4paper,UTF8]{ctexart}
\usepackage[a4paper,top=22mm,bottom=22mm,left=22mm,right=22mm,headheight=15pt]{geometry}
\usepackage{amsmath,amssymb,mathtools,bm}
\usepackage{graphicx,xcolor,booktabs,longtable,array,enumitem,fancyhdr,caption,float,microtype,setspace}
\newcommand{\Needspace}[1]{\par\begingroup\dimen0=#1\relax\dimen2=\pagegoal\advance\dimen2 by-\pagetotal\ifdim\dimen0>\dimen2\newpage\fi\endgroup}
\newcommand{\needspace}[1]{\par\begingroup\skip0=\lastskip\vskip-\skip0\dimen0=#1\relax\advance\dimen0 by\skip0\vskip\dimen0\penalty-100\vskip-\dimen0\vskip\skip0\penalty9999\endgroup}
\usepackage[hidelinks,unicode]{hyperref}
\usepackage[most]{tcolorbox}
\definecolor{ContestBlue}{HTML}{17365D}
\hypersetup{pdftitle={第 046 套},pdfauthor={}}
\setstretch{1.24}
\setlength{\parindent}{2em}
\setlength{\parskip}{0.35em}
\setlength{\emergencystretch}{2em}
\setlist{itemsep=0.35em,topsep=0.4em}
\captionsetup{font=small,labelfont=bf,justification=centering}
\setcounter{secnumdepth}{0}
\pagestyle{fancy}
\fancyhf{}
\fancyhead[L]{\small 第 046 套}
\fancyfoot[C]{\thepage}
\renewcommand{\headrulewidth}{0.4pt}
\ctexset{
  section={format=\Large\bfseries\color{ContestBlue},beforeskip=1.2em,afterskip=0.7em},
  subsection={format=\large\bfseries\color{ContestBlue},beforeskip=1.0em,afterskip=0.5em},
  subsubsection={format=\normalsize\bfseries,beforeskip=0.8em,afterskip=0.35em}
}
\makeatletter
\newcommand{\tightlist}{\setlength{\itemsep}{0pt}\setlength{\parskip}{0pt}}
\newcommand{\maxwidth}{\ifdim\Gin@nat@width>\linewidth\linewidth\else\Gin@nat@width\fi}
\newcommand{\maxheight}{\ifdim\Gin@nat@height>0.72\textheight 0.72\textheight\else\Gin@nat@height\fi}
\newcommand{\pandocbounded}[1]{\begingroup\setkeys{Gin}{width=\maxwidth,height=\maxheight,keepaspectratio}#1\endgroup}
\makeatother
\begin{document}
\begin{titlepage}
  \centering
  \vspace*{58mm}
  {\Huge\bfseries\color{ContestBlue} 第 046 套\par}
  \vfill
  \vspace*{28mm}
\end{titlepage}
\section{第 1 题}\label{ux7b2c-1-ux9898}

地球在某种假想作用下变成无限长、均匀的质量圆柱，仍保持半径
\(R=6400\,\mathrm{km}\) 和密度
\(\rho=5.5\times10^3\,\mathrm{kg\,m^{-3}}\)。月球仍可视为球形质点，只考虑地球---月球系统。

（1）在地球改变之前，已知月球公转周期为 \(T=27.3\) 天、月球质量为
\(7.35\times10^{22}\,\mathrm{kg}\)，按匀速圆周运动求地月距离和月球速度。

（2）地球改变为无限长质量圆柱后，若月球仍能绕圆柱轴线做匀速圆周运动，求速度
\(v\) 与月球到轴线距离 \(r\) 的关系。

\begin{figure}[H]
\centering
\includegraphics[width=0.8\linewidth,height=\textheight,keepaspectratio,alt={题1　球形地球与无限长质量圆柱两种引力模型}]{../figures/question-01.pdf}
\caption{题1　球形地球与无限长质量圆柱两种引力模型}
\end{figure}

\newpage

\section{第 2 题}\label{ux7b2c-2-ux9898}

一辆车在水平地面上加速。车厢角落放有一个半径为 \(r\)
的圆形刚体，其重心位于几何圆心正下方 \(h\) 处。地面摩擦系数为
\(\mu\)，左侧墙壁光滑，重力加速度为 \(g\)。

（1）车向左或向右加速时，哪种情形可能使圆形刚体发生滑动？选项为：A.
左；B. 右；C. 都可能；D. 都不可能。

（2）若车向右加速，发生滑动所需的最小加速度是多少？

（3）若地面光滑、墙壁摩擦系数为
\(\mu'\)，求保证无论加速度多大都不滑动所需的 \(\mu'\) 最小值。

\begin{figure}[H]
\centering
\includegraphics[width=0.702\linewidth,height=\textheight,keepaspectratio,alt={题2　题设与符号}]{../figures/question-02.pdf}
\caption{题2　题设与符号}
\end{figure}

\newpage

\section{第 3 题}\label{ux7b2c-3-ux9898}

两根无限长直导线垂直于 \(xOy\) 平面，分别通过 \((a,0)\) 和
\((-a,0)\)，单位长度带电量为 \(-\lambda\) 和
\(+\lambda\)。由此建立电场线与薄圆盘稳恒电流的类比。

（1）\(y\) 轴上点 \((0,\eta a)\) 的场强写成 \[
E_x(\eta)=\frac{\zeta\lambda}{4\pi\varepsilon_0a},
\] 求 \(\zeta\)。

（2）通过 \((0,\pm\eta a)\) 的两条电场线在 \(xOy\) 平面围成面积
\(S=\xi a^2\)，求 \(\xi\)。

（3）半径为 \(r\)、厚度 \(h\ll r\)、电导率为 \(\sigma_0\)
的均匀圆盘，在直径两端接入半径 \(\rho\ll r\)
的圆柱形电极，左正右负。圆心电流密度为 \(j_0\)，边缘极角 \(\theta\)
处电势为 \(U(\theta)\)。若 \[
\left|U\!\left(\frac\pi2\right)-U(\theta)\right|
=\gamma\frac{j_0r}{\sigma_0},
\] 求 \(\gamma\)。

（4）保持两极电压不变，把两电极移到圆盘边缘、角距离为 \(\phi\)
的两点，圆心电流密度变为 \(\chi j_0\)，求 \(\chi\)。

（5）保持第（3）问电压，把电极改到
\((\pm\mu r,0)\)，\(0<\mu<1\)，圆心电流密度变为 \(\chi'j_0\)，求
\(\chi'\)。

\begin{figure}[H]
\centering
\includegraphics[width=0.8\linewidth,height=\textheight,keepaspectratio,alt={题3　异号线电荷的电场线与圆盘电极几何}]{../figures/question-03.pdf}
\caption{题3　异号线电荷的电场线与圆盘电极几何}
\end{figure}

\newpage

\section{第 4 题}\label{ux7b2c-4-ux9898}

一根直径为 \(D\) 的玻璃管弯成 U 形，高为 \(h\)、宽为
\(3l\)，内部装有密度为 \(\rho\) 的理想液体。平衡时左右液柱高度均为
\(h/2\)，重力加速度为 \(g\)。

（1）求液面小幅上下振动的角频率 \(\omega_1\)。

（2）U 形管绕一条竖直轴以恒定角速度转动，转轴距左侧竖直管
\(l\)；平衡时左液面下降 \(\Delta h\)，求角速度 \(\omega_0\)。

（3）求第（2）问恒定角速度转动时液面小振动角频率 \(\omega_2\)。

（4）初态以 \(\omega_0\)
转动，随后撤去外界对角速度的限制，系统依靠惯性继续转动，求液面小振动角频率
\(\omega_3\)。

\begin{figure}[H]
\centering
\includegraphics[width=0.726\linewidth,height=\textheight,keepaspectratio,alt={题4　U 形液柱及偏置竖直转轴}]{../figures/question-04.pdf}
\caption{题4　U 形液柱及偏置竖直转轴}
\end{figure}

\newpage

\section{第 5 题}\label{ux7b2c-5-ux9898}

水平桌面上有一条光滑的抛物线形导轨，方程为 \(y=kx^2\)，其中 \(k\)
为正常数，\(xOy\) 平面为水平面。一根与 \(x\)
轴平行的导体棒跨接在导轨上，两端与导轨保持接触。空间中存在方向竖直、磁感应强度大小为
\(B\) 的匀强磁场。导体棒从坐标原点开始沿 \(y\) 轴正方向以恒定速率 \(v\)
运动。

（1）求回路感应电动势 \(E(t)\)。

（2）导体棒的单位长度电阻为 \(\lambda\)，导轨电阻不计，求回路电流。

（3）求维持导体棒匀速运动所需外力的功率。

\begin{figure}[H]
\centering
\includegraphics[width=0.8\linewidth,height=\textheight,keepaspectratio,alt={题5　抛物线导轨与匀速横杆}]{../figures/question-05.pdf}
\caption{题5　抛物线导轨与匀速横杆}
\end{figure}

\newpage

\section{第 6 题}\label{ux7b2c-6-ux9898}

余弦发光体是指：面元 \(dS\) 在时间 \(dt\) 内向立体角 \(d\Omega\)
发出的能量为 \[
dE=I(\theta)d\Omega\,dS\,dt,
\qquad I(\theta)=I_0\cos\theta,
\] 其中 \(\theta\) 为出射方向与面元法线的夹角。亮度定义为 \[
B=\frac{dE}{d\Omega\,dS\,dt\cos\theta}.
\]

（1）证明均匀球形余弦发光体从远处看亮度均匀。

（2）均匀余弦发光平面上放置半径为 \(r\)、折射率为 \(n=2\)
的玻璃半球，半球平面与发光平面紧密贴合。假设未发生全反射的光线均能透射。观察者位于发光平面法线方向的远处，求玻璃半球投影区域内的亮度分布，忽略经历两次反射的光线。

\begin{figure}[H]
\centering
\includegraphics[width=0.86\linewidth,height=\textheight,keepaspectratio,alt={题6　题设与符号}]{../figures/question-06.pdf}
\caption{题6　题设与符号}
\end{figure}

\newpage

\section{第 7 题}\label{ux7b2c-7-ux9898}

一个球形小液滴悬浮在空气中，折射率为 \(n\)。一束平行光以入射角
\(\alpha\) 入射，经液滴内部 \(N\) 次反射后出射。定义光线总偏转角为
\(\Delta_N\)。若 \(\Delta_N(\alpha)\)
出现极值，则附近一簇入射光会集中到相近方向，形成第 \(N\) 级虹；\(N=1\)
称为虹，\(N=2\) 称为霓。

（1）给定 \(n=n_0\)、\(\alpha=\alpha_0\)，求 \(\Delta_1(\alpha_0)\)。

（2）求 \(\Delta_2(\alpha_0)\)。

（3）对任意 \(n,N\)，求极值对应的入射角 \(\alpha_{Nm}\)。

（4）求极值 \(\Delta_{Nm}\)。

（5）为使第 \(N\) 级虹出现，求 \(n\) 可取的最大值 \(n_{Nm}\)。

（6）已知一种偏振分量从空气入射液滴时的振幅反射率为 \[
r=\frac{\cos\alpha-n\cos\theta}
{\cos\alpha+n\cos\theta},
\] 透射率为 \(t\)；从液滴入射空气时 \(r'=-r\)、透射率为 \(t'\)，且
\(tt'=1+rr'\)。取 \(n_0=2.03\)，求第 2 级虹与第 1 级虹的光振幅平方之比
\(I_2/I_1\)。

（7）求第 3 级虹与第 1 级虹的光振幅平方之比 \(I_3/I_1\)。

\begin{figure}[H]
\centering
\includegraphics[width=0.43\linewidth,height=\textheight,keepaspectratio,alt={题7　液滴内一次反射光路与偏转角}]{../figures/question-07.pdf}
\caption{题7　液滴内一次反射光路与偏转角}
\end{figure}

\newpage

\section{第 8 题}\label{ux7b2c-8-ux9898}

设飞船可控地进行质子聚变： \[
4\mathrm H^+\longrightarrow
\mathrm{He}^{2+}+2e^++2\nu_e+n_\gamma\gamma.
\] 已知 \[
\begin{aligned}
m_{\mathrm H}&=1.672622\times10^{-27}\ \mathrm{kg},\\
m_{\mathrm{He}}&=6.644657\times10^{-27}\ \mathrm{kg},\\
m_e&=9.109384\times10^{-31}\ \mathrm{kg},\\
c&=299792458\ \mathrm{m\,s^{-1}}.
\end{aligned}
\]

（1）为最有效地加速飞船，应把反应释放的能量集中在哪一种产物上并反向喷出？

（2）若全部反应能量集中到该粒子，求其相对飞船的速度 \(u\)。

（3）质子反应速率为 \(1\ \mathrm{g\,s^{-1}}\)，求飞船参考系中的最大推力
\(F\)。

（4）地面参考系中，飞船速度为 \(v\)，飞船及燃料的瞬时静质量为
\(m_0\)，求地面观察到的最大加速度 \(a'\)。

（5）把飞船从静止加速到 \(0.20c\)，求初始燃料质量与飞船本身质量之比
\(\eta\) 的最小值。

（6）飞船最多携带相当于飞船本身质量 \(100\)
倍的燃料。一次航行需在自身参考系中至多巡航 \(15\)
年，并在下一恒星处近乎静止以补充燃料和维护。忽略加速、减速阶段所占时间，求一次航行最大行程
\(S\)，以光年表示。

\begin{figure}[H]
\centering
\includegraphics[width=0.8\linewidth,height=\textheight,keepaspectratio,alt={题8　聚变反冲推进与星际航行阶段}]{../figures/question-08.pdf}
\caption{题8　聚变反冲推进与星际航行阶段}
\end{figure}
\end{document}
